% ════════════════════════════════════════════════════════════════════
% Chapter 9 — Validation and Quality Assurance
% ════════════════════════════════════════════════════════════════════
\section{Validation and Quality Assurance}\label{sec:validation}

A research-grade data collection platform must provide verifiable evidence that every stage of the measurement chain---from the ADC conversion to the final kinematic trajectory---meets quantitative accuracy targets.
This section describes the multi-tiered validation framework, comprising pre-session hardware verification, in-session real-time monitoring, and post-session signal-processing validation.


\subsection{Pre-Session Validation: Pre-Flight Checklist}

Before any recording can begin, the GUI enforces a systematic pre-flight checklist (\class{PreflightPanel}).
The START button is gated; only when all checks pass (or the operator explicitly overrides with documented justification) does the system transition to the RECORDING state.

\begin{table}[H]
  \centering
  \caption{Pre-flight checklist items and acceptance criteria.}
  \label{tab:preflight}
  \begin{tabular}{p{3.5cm}p{3.5cm}p{3.5cm}p{3cm}}
    \toprule
    \textbf{Check} & \textbf{Pass} & \textbf{Warn} & \textbf{Fail} \\
    \midrule
    IMU sample rate     & Within $\pm 5\%$ of \SI{1}{\kilo\hertz} (\SI{950}{}--\SI{1050}{\hertz}) & Within $\pm 10\%$ & Outside $\pm 10\%$ or no samples \\
    Accel saturation    & 0 clipped samples  & $< 10$ clipped & $\geq 10$ clipped \\
    Camera connected    & \code{is\_open() == true}  & --- & Not connected \\
    Marker detection    & Detection rate $> 80\%$ & $> 50\%$ & $\leq 50\%$ \\
    Marker SNR          & $> 3.0$              & $> 2.0$ & $\leq 2.0$ \\
    Sync drift          & $< \SI{50}{ppm}$     & $< \SI{100}{ppm}$ & $\geq \SI{100}{ppm}$ or rearm flagged \\
    Subject ID          & Non-empty            & --- & Empty \\
    \bottomrule
  \end{tabular}
\end{table}

Any override of a failed check is recorded as a \code{preflight\_override} event in \file{events.jsonl}, including the specific check that was overridden and the operator's reason, preserving the audit trail for downstream data quality review.


\subsection{In-Session Monitoring}

\subsubsection{IMU Signal Quality Indicators}

The following metrics are computed in real-time and displayed in the Sensor panel:

\begin{itemize}[nosep]
  \item \textbf{Measured rate:} Rolling estimate of the actual IMU sample rate from ESP32 timestamps. Deviations $> 5\%$ from nominal trigger a toast warning.
  \item \textbf{Jitter statistics:} Mean, max, and standard deviation of inter-sample timing jitter. Jitter $> \SI{100}{\micro\second}$ indicates potential SPI bus contention or RTOS priority inversion.
  \item \textbf{Per-axis saturation count:} Cumulative count of samples at $\pm$32768 (full-scale ADC). Persistent saturation indicates an insufficient FSR.
  \item \textbf{Gyro noise floor:} Rolling standard deviation of $|\boldsymbol{\omega}|$ during detected rest periods. Expected $< \SI{0.1}{\degree\per\second}$ for the ICM-42688-P.
  \item \textbf{Accel noise floor:} Rolling standard deviation of $|\mathbf{a}| - 1$ at rest. Expected $< \SI{0.005}{\g}$.
  \item \textbf{CRC error rate:} Cumulative count and rate of packets failing the CRC16-CCITT check.
  \item \textbf{Dropout count:} Number of gaps in the ESP32 timestamp sequence exceeding \SI{2}{\milli\second}.
\end{itemize}

\subsubsection{FSYNC Health Monitoring}

The FSYNC hit rate is tracked in real-time and should match the camera frame rate ($\sim$\SI{90}{\hertz}).
A mismatch indicates a fault in the level-shifting circuit, a loose SYNC cable, or a camera configuration error.

\subsubsection{Synchronisation Drift Monitoring}

The \class{SyncEngine} continuously computes the clock drift between the ESP32 and D455 hardware clocks using the FSYNC anchor pairs (\cref{sec:sync_engine}).
If drift exceeds \SI{100}{ppm} for \SI{5}{\second}, the auto-rearm mechanism is triggered, halting data collection until the operator performs a fresh tap test.

\subsubsection{Per-Repetition Kinematic Plausibility}

After each repetition is segmented, the \class{Validator} applies configurable plausibility bounds:

\begin{table}[H]
  \centering
  \caption{Per-repetition kinematic plausibility bounds (configurable per exercise profile).}
  \label{tab:plausibility}
  \begin{tabular}{llcc}
    \toprule
    \textbf{Metric} & \textbf{Unit} & \textbf{Min} & \textbf{Max} \\
    \midrule
    Mean Concentric Velocity & \si{\metre\per\second} & 0.05 & 3.00 \\
    Peak Concentric Velocity & \si{\metre\per\second} & 0.10 & 5.00 \\
    Range of Motion          & \si{\metre}            & 0.05 & 1.50 \\
    Concentric Duration      & \si{\second}           & 0.10 & 10.00 \\
    \bottomrule
  \end{tabular}
\end{table}

Repetitions that violate any bound are flagged as implausible.
The Rep Timeline Panel displays a colour-coded dot for each rep: green for passed, red for failed.
Hovering over a red dot shows the specific violations (e.g., ``PCV exceeds \SI{5.00}{\metre\per\second}: likely double-counted bounce'').


\subsection{Post-Session Validation}\label{sec:post_session_validation}

% The "Golden Model Pipeline" that used to be described here has been removed from the
% project; the camera--IMU agreement figures it reported are superseded by
% Section~\ref{sec:imu_baseline}, where every number is produced by a named script.

Post-session validation has two parts, and they are separated deliberately because they
answer different questions.

\paragraph{Is the reference sound?} The smoother's own consistency is checked by its
normalised innovation squared, which is driven to approximately unity by the choice of
process noise rather than assumed --- see Section~\ref{sec:sync_measured} for the frame
grid this runs on. The resulting per-frame uncertainties are the reference's error bars and
are released with the corpus: \SI{0.69}{\milli\metre} of height,
\SI{8.1}{\milli\metre\per\second} of speed. A separate check confirms the reference does
not flatter itself: the residual in its velocity above \SI{6}{\hertz} is
\SI{0.08}{\milli\metre\per\second}, and the peak it reports is unchanged whether picked
from the smoother output or from a further-smoothed version, so it contributes nothing
measurable to the peak-velocity comparison.

\paragraph{Does the inertial estimate agree with it?} That is the subject of
Section~\ref{sec:imu_baseline}, which reports agreement per repetition with bias, limits of
agreement, standard error of the estimate and correlation, both given the reference's
repetition boundaries and with no camera input at all.

\subsection{Synchronisation Validation: Tap Test}

The tap test (\cref{sec:sync_engine}) serves as both calibration and validation.
After the test, the system reports:
\begin{itemize}[nosep]
  \item \textbf{Temporal offset:} The systematic IMU-to-camera delay. Expected $< \SI{1}{\milli\second}$ with the hardware FSYNC path.
  \item \textbf{Drift rate:} The rate of clock divergence. Expected $< \SI{50}{ppm}$ (\SI{0.18}{\second\per\hour}).
  \item \textbf{Cross-correlation quality:} The normalised cross-correlation coefficient between the IMU acceleration magnitude and the camera position displacement. Expected $> 0.9$.
\end{itemize}

Multiple tap tests within a session can be used to verify that drift remains bounded throughout the recording duration.


\subsection{Allan Variance Characterisation}

The sensor is characterised per IEEE~952-2020~\cite{ieee952}, and \emph{from the corpus
itself} rather than from a dedicated static recording. The overlapping Allan deviation
$\sigma(\tau)$ is computed over \SI{507}{\second} of stillness in 72 stretches, found in the
sensor's own angular rate across all \num{84} sessions. Stillness is \emph{found}, never
assumed from a position in the session: the interval before the first repetition is not
still, because the bar is being handled, and reading noise off it gives figures two hundred
times the datasheet.

\begin{table}[t]
  \centering
  \caption{Sensor characterisation, from \SI{507}{\second} of detected stillness across the
           corpus. Produced by \texttt{scripts/imu/noise\_characterisation.py}.}
  \label{tab:allan}
  \begin{tabular}{lr}
    \toprule
    \multicolumn{2}{l}{\emph{Gyroscope}} \\
    White noise, angle random walk & \SI{1.93}{\milli\degree\per\second\per\sqrt\hertz} \\
    Bias instability               & \SI{12.2}{\degree\per\hour} \\
    Rate random walk               & \SI{7.67}{\milli\degree\per\second\squared\per\sqrt\hertz} \\
    \midrule
    \multicolumn{2}{l}{\emph{Accelerometer}} \\
    White noise, velocity random walk & \SI{50.1}{\micro\gram\per\sqrt\hertz} \\
    Bias instability                  & \SI{100}{\micro\gram} \\
    \bottomrule
  \end{tabular}
\end{table}

Two consequences. First, these are the values the process-noise covariance
$\mathbf{Q}$ actually needs, and they matter: the filter had previously been run with a
gyroscope white-noise density set by hand at twenty-six times the measured figure, which
makes it distrust the gyroscope and lean on an accelerometer being shaken by the lift.
Second, they bound what the sensor can support. Over a \SI{1.5}{\second} concentric these
densities permit roughly \SI{20}{} to \SI{25}{\milli\metre\per\second}, against the
\SI{47.9}{\milli\metre\per\second} measured in Section~\ref{sec:imu_baseline} --- so the
sensor is not the limiting term.


\subsection{Diagnostic Export}

At any time during or after a session, the operator can trigger a diagnostic bundle export (F5 or Tools $\rightarrow$ Export Diagnostics).
This produces a PII-free JSON file containing:
\begin{itemize}[nosep]
  \item Build provenance (version, git SHA, compiler).
  \item Current IMU/camera/sync statistics.
  \item Recent toast notification history.
  \item Full application configuration (with PII fields redacted).
\end{itemize}

This bundle enables remote troubleshooting without requiring access to the raw session data.


\subsection{Validation Summary and Acceptance Criteria}

\begin{table}[H]
  \centering
  \caption{System-level validation targets.}
  \label{tab:validation_summary}
  \begin{tabular}{lcc}
    \toprule
    \textbf{Metric} & \textbf{Target} & \textbf{Measurement Method} \\
    \midrule
    Position RMSE (ESKF vs.\ camera)    & $< \SI{50}{\milli\metre}$         & Golden model Step~8 \\
    Velocity bias                       & $< \SI{0.05}{\metre\per\second}$   & Bland--Altman mean difference \\
    Velocity LoA (95\%)                 & $< \pm\SI{0.10}{\metre\per\second}$ & Bland--Altman 1.96$\sigma$ \\
    Synchronisation offset              & $< \SI{1}{\milli\second}$         & Tap test \\
    Clock drift                         & $< \SI{50}{ppm}$                  & FSYNC anchor pairs \\
    IMU sample rate stability           & $\pm 5\%$ of \SI{1}{\kilo\hertz}  & ESP32 timestamp analysis \\
    CRC error rate                      & $< 0.01\%$                        & Packet validation \\
    Marker detection rate               & $> 80\%$                          & Session statistics \\
    \bottomrule
  \end{tabular}
\end{table}
